% sec0.tex — c01｜原文 p2｜§0 Plan of lectures（讲义提纲）
% 0 条编号公式；无编号环境。
% [原文 p.2]
超共形指标是超共形场论在径向量子化下的 Witten 指标。我们将首先在第 1 节结合超对称量子力学介绍 Witten 指标。除了理解指标受保护背后的基本原理，我们还将发展许多工具，特别是齐次指数化（plethystic exponentiation）；这些工具对于在更简单的设定下计算规范理论的超共形指标是必需的。在第 2 节，我们考虑超共形代数并引入超共形指标。基本组元——即手征多重态与矢量多重态——的指标将被计算出来。随后我们讨论指标在超共形代数短多重态重组下的不变性。我们的结论是：超共形指标在重组的意义下（up to recombinations）捕捉了理论的全部谱信息。在第 3 节，我们把超共形指标应用于所谓的「class S」对偶。在解释了具有超对称性增强的指标极限之后，我们描述与指标相伴的拓扑场论。最后，在第 4 节，我们计算大 $N$ 极限下 $\sN=4$ 的 $\SU(N)$ 超 Yang–Mills 的指标。先验地，$AdS_5$ 中的超对称黑洞被认为会对指标贡献 $O(N^2)$ 个态。
% [原文 p.3]
但我们发现，指标是 $O(1)$ 的，并且只由体（bulk）中的超引力态所饱和。我们对这一表面上的不匹配加以评论，并介绍最近一些从指标恢复黑洞熵的工作。指标的其它一些方面，例如它的模性质 [1, 2] 以及它与反常的联系 [3]，将不在这些讲义中讨论。我们只讨论 $S^3$ 上的指标。关于透镜空间（lens space）上指标的讨论，见 [4–7]。关于超共形指标的其它一些优秀综述，见 [8, 9]。

% [Tutor Note: 「径向量子化」把 $S^{d-1}\times\R$ 上的理论化为以径向为时间的哈密顿框架；超共形指标正是该框架下的 Witten 指标。]
% [Tutor Note: 「受保护」（protected）指该量在保持超对称的形变下不变，因而可在弱耦合自由极限精确计算。]
% [Tutor Note: 齐次指数把单字母谱提升为多粒子谱；它是第 1 节的技术核心，后文 (2.6)、(2.13) 会反复使用。]
% [Tutor Note: 「重组」（recombination）指短多重态合并成长多重态；指标对重组不变，故只捕捉「重组意义下」的谱信息。]
% [Tutor Note: 原文 class S 带引号且用花体 S；glossary 冻结为 class S 原样保留，可称「S 类理论」。]
% [Tutor Note: $AdS_5$ 黑洞的 $O(N^2)$ 微观态与指标 $O(1)$ 的落差是 §4 的主题；此处先给结论，细节见 (4.16) 附近。]
% [Tutor Note: 「透镜空间」（lens space）上的指标涉及更一般的边界条件；本讲义只讨论 $S^3$。]
% [Tutor Note: 本节是全篇路线图：可先跳到 §1 的谐振子例子建立直觉，再回看 §2 的公理化定义。]
